\chapter{Atividades Futuras}
\label{cap:futuros}

Concluídas a implementação da ferramenta e a sua avaliação de desempenho, esta seção
reúne as atividades que dão continuidade ao trabalho. Elas dividem-se entre a ampliação da
forma de uso, o aprofundamento da avaliação e a extensão das funcionalidades de importação.

\begin{itemize}
  \item \textbf{Interface gráfica \textit{desktop}.} A ferramenta é hoje operada por uma
        interface de linha de comando (Seção~\ref{sec:visao-geral}). O próximo passo
        desejável --- e viável ainda dentro deste ciclo de trabalho --- é uma interface
        gráfica que monte o comando de importação de forma assistida e permita visualizar e
        editar os dois arquivos JSON de configuração, reduzindo a barreira de entrada para
        usuários menos familiarizados com a linha de comando.

  \item \textbf{Aprofundamento da avaliação de desempenho.} A avaliação do
        Capítulo~\ref{cap:avaliacao} cobriu volumes de até cem mil linhas em uma única
        máquina, com os cinco SGBDs em contêineres partilhando os mesmos recursos. Três
        extensões se destacam: medir volumes maiores, na ordem de milhões de linhas;
        aprofundar o caminho de escrita do Cassandra --- a fase dominante do custo --- para
        além da concorrência já empregada pela estratégia otimizada
        (Seção~\ref{sec:estrategias-escrita}); e repetir os experimentos em um ambiente com
        SGBDs dedicados e ajustados, isolando o desempenho da ferramenta das restrições do
        ambiente de contêineres.

  \item \textbf{Suporte a outros formatos de entrada.} Aceitar arquivos TSV e planilhas
        Excel, além de CSV, reaproveitando o mesmo mapeamento declarativo já existente.

  \item \textbf{Operadores de filtro adicionais e políticas de coerção configuráveis.}
        Ampliar o vocabulário de seleção de linhas e permitir que o usuário defina como
        valores ambíguos ou malformados são convertidos aos tipos nativos de cada SGBD.

  \item \textbf{Testes de integração em ambiente conteinerizado.} A suíte de testes atual é
        majoritariamente unitária; prevê-se automatizar testes de integração contra os cinco
        SGBDs em \texttt{docker compose}, dentro de um fluxo de integração contínua (CI).

  \item \textbf{Generalização do núcleo de importação.} Estender a arquitetura de portas e
        adaptadores (Seção~\ref{sec:modelos-execucao}) a novos destinos --- sistemas de
        mensageria, \textit{data lakes} e outros SGBDs ---, reaproveitando o orquestrador
        agnóstico de destino sem duplicar a lógica de importação.
\end{itemize}
